Put \(q=m-2c=h-c\), which is at least three by the theorem's stated range. The scale bound implies \(0\leq s_{ek}\leq\overline M-1\). For any fixed coordinate pair, the absolute affine determinant of three points in the square \([0,\overline M-1]^2\) is at most \((\overline M-1)^2\): holding five of its six scalar coordinates fixed, it is the absolute value of an affine function of the sixth coordinate, so its maximum over that coordinate interval occurs at an endpoint; iterating over all six coordinates reduces the maximum to triples of square vertices, whose determinants have absolute value zero or \((\overline M-1)^2\). Hence \(\gamma_H\leq(\overline M-1)^2\). The declared restrictions \(q\geq3\), \(0<\gamma_0\leq(\overline M-1)^2\), and \(0<\zeta_0\leq1\) therefore explicitly remove the elementary vacuous parameter regimes.

\textbf{Uniform structural bounds.}
Let \(G\) be either the true structural matrix or any feasible candidate structural matrix satisfying the stated cycle-slack and condition-number bounds. Since \(\operatorname{diag}(G)=\mathbf 1\), the identity-permutation term in the Leibniz expansion of \(\det G\) is one. Each nonidentity permutation is a product of fixed points and pairwise disjoint directed simple cycles of length at least two. For an off-diagonal entry, \((I_p-G)_{ij}=-G_{ij}\). Hence the absolute product contributed by each nontrivial cycle is at most
\[
\operatorname{CP}(I_p-G)\leq1-\zeta_0<1.
\]
The assumption \(0<\zeta_0\leq1\) makes the right-hand side an element of \([0,1)\). It follows that the absolute value of every one of the \(p!\) Leibniz terms is at most one, and therefore
\[
|\det G|\leq p!.
\tag{1}
\]
Order the singular values as \(\sigma_1(G)\geq\cdots\geq\sigma_p(G)>0\). The condition-number bound gives \(\sigma_j(G)\geq\sigma_p(G)\geq\sigma_1(G)/\overline\kappa\) for \(2\leq j\leq p\). Thus
\[
\frac{\sigma_1(G)^p}{\overline\kappa^{p-1}}
\leq\prod_{j=1}^p\sigma_j(G)
=|\det G|
\leq p!,
\]
so
\[
\|G\|_{\mathrm{op}}\leq L_0.
\tag{2}
\]
Moreover, a unit diagonal entry implies \(\|G\|_{\mathrm{op}}\geq1\). Consequently,
\[
\|G^{-1}\|_{\mathrm{op}}
=\frac{\|G\|_{\mathrm{op}}\|G^{-1}\|_{\mathrm{op}}}
{\|G\|_{\mathrm{op}}}
\leq\overline\kappa.
\tag{3}
\]

\textbf{A compact residual gap without covariance interiority.}
Set
\[
\rho_0=\frac{1}{2K_0\overline M}>0.
\tag{4}
\]
For every choice of \(H_1,H_2\subseteq E\) with \(|H_1|=|H_2|=h\) and every \(S\subseteq H_1\cap H_2\) with \(|S|=q\), consider the finite-dimensional tuples
\[
\tau=\bigl(D,U,\Omega,(s_e,\Sigma_e)_{e\in H_1},
A,V,\Psi,(t_e,\Gamma_e)_{e\in H_2}\bigr)
\]
that satisfy all of the following conditions. First, \(DU=UD=I_p\), \(AV=VA=I_p\), and both \(D\) and \(A\) have unit diagonal. Second,
\[
\operatorname{CP}(I_p-D)\leq1-\zeta_0,
\quad
\operatorname{CP}(I_p-A)\leq1-\zeta_0,
\quad
\|D\|_{\mathrm{op}}\|U\|_{\mathrm{op}}\leq\overline\kappa,
\quad
\|A\|_{\mathrm{op}}\|V\|_{\mathrm{op}}\leq\overline\kappa.
\]
Third, \(\Omega,\Psi,\Sigma_e,\Gamma_e\succeq0\), the vectors \(s_e,t_e\) are coordinatewise nonnegative, and each side obeys its BACKSHIFT covariance equations. Fourth, each side has affine separation at least \(\gamma_0\), computed over every \(q\)-element subset of its size-\(h\) index set. Finally, each side obeys the corresponding scale bound with upper bound \(\overline M\). Let \(\mathfrak T\) be the union of these tuple sets over the finitely many choices of \(H_1,H_2,S\).

Equations (2)--(3) bound \(D,U,A,V\). Because all summands in each scale inequality are nonnegative, the scale inequalities separately bound \(\Omega,\Psi\), every covariance matrix, and every coordinate of every nonnegative shift vector. All defining constraints are closed: the inverse and covariance relations are polynomial equalities; the positive-semidefinite cones and nonnegative orthants are closed; matrix norms are continuous; \(\operatorname{CP}\) is the maximum of finitely many absolute monomials; and each affine-separation constraint is obtained from finitely many maxima and minima of absolute determinants. Thus every fixed-index tuple set is closed and bounded in a finite-dimensional Euclidean space. By \(\text{lem:finite-dimensional-compact-extrema}\), it is compact, and the finite union \(\mathfrak T\) is compact. Notice that this argument uses only positive semidefiniteness and the scale bound; it nowhere uses a lower covariance eigenvalue.

Define the continuous residual
\[
\mathfrak r(\tau)=
\max_{e\in S}
\left\|A\Sigma_eA^{\mathsf T}
-\Psi-\operatorname{diag}(t_e)\right\|_{\mathrm F}
\]
and the closed, hence compact, far family
\[
\mathfrak T_{\mathrm{far}}
=\{\tau\in\mathfrak T:\|AU-I_p\|_{\mathrm F}\geq\rho_0\}.
\]
If \(\mathfrak T_{\mathrm{far}}=\varnothing\), set \(\mu_0=1\). Otherwise \(\text{lem:finite-dimensional-compact-extrema}\) gives an attained minimum
\[
\mu_0=\min_{\tau\in\mathfrak T_{\mathrm{far}}}\mathfrak r(\tau).
\tag{5}
\]
This minimum is strictly positive. Indeed, if a far tuple had \(\mathfrak r(\tau)=0\), then for every \(e\in S\),
\[
D\Sigma_eD^{\mathsf T}=\Omega+\operatorname{diag}(s_e),
\qquad
A\Sigma_eA^{\mathsf T}=\Psi+\operatorname{diag}(t_e).
\tag{6}
\]
For each \(k<\ell\), the first side's affine-separation condition, applied to the particular \(q\)-set \(S\), supplies three labels in \(S\) whose \((k,\ell)\) affine minor has absolute value at least \(\gamma_0>0\). Hence the first side's \((k,\ell)\)-shift cloud on \(S\) is not contained in an affine line.

Since \(|E\setminus S|=m-q=2c\), choose a disjoint partition \(E\setminus S=C_1\mathbin{\dot\cup}C_2\) with \(|C_1|=|C_2|=c\). Define one covariance family on \(E\) by
\[
\widetilde\Sigma_e=
\begin{cases}
\Sigma_e,&e\in S,\\
D^{-1}\Omega D^{-\mathsf T},&e\in C_1,\\
A^{-1}\Psi A^{-\mathsf T},&e\in C_2.
\end{cases}
\]
The first explanation fits \(S\cup C_1\) using shifts \(s_e\) on \(S\) and zero shifts on \(C_1\); the second fits \(S\cup C_2\) using shifts \(t_e\) on \(S\) and zero shifts on \(C_2\). Both newly assigned covariances are positive semidefinite by congruence of \(\Omega,\Psi\succeq0\), and both fitted sets have size \(q+c=h\) and intersection \(S\). Equation (6), nonnegativity, the normalization constraints, and the established noncollinearity verify every hypothesis of \(\text{lem:overlap-uniqueness}\). That lemma gives \(A=D\), whence \(AU=I_p\), contradicting the far-family inequality. Therefore
\[
\mu_0>0.
\tag{7}
\]
This also holds under the convention \(\mu_0=1\) when the far family is empty. Define, non-effectively,
\[
r_0=\frac{\mu_0}{2\sqrt p\,L_0^2}>0.
\tag{8}
\]
The compact family defining \(\mu_0\) uses only \(p,m,c,\gamma_0,\zeta_0,\overline\kappa,\overline M\), so \(r_0\) has precisely the claimed dependence. This definition neither decides whether the far family is empty nor computes its minimum.

\textbf{Entry into the local chart.}
Fix an outcome in \(G_{n,\alpha,H}\), write \(r=r_{n,\alpha,H}\), assume \(r\leq r_0\), and take \(A\in\mathcal U^{(n)}_\alpha\). By \(\text{def:confidence-union}\), there is a simultaneous selection \(\Gamma_e\in\mathcal C^{(n)}_{e,\alpha}\), \(e\in E\), for which \(A\in\mathcal R_c(\Gamma)\). By \(\text{def:compatible-set}\), a set of at least \(h\) labels, a nuisance \(\Psi\succeq0\), and nonnegative shifts satisfy the candidate covariance equations. Restrict these witnesses to a set \(K\subseteq E\) of exactly \(h\) labels. The declarative premise supplies all four candidate-side bounds for this restricted tuple.

Inclusion--exclusion yields
\[
|H\cap K|\geq|H|+|K|-|E|=2h-m=q.
\]
Choose \(S\subseteq H\cap K\) with \(|S|=q\), and put \(U=D^{-1}\) and \(V=A^{-1}\). The true and candidate equations and four-margin bounds place the resulting objects in \(\mathfrak T\). Since the outcome is in \(G_{n,\alpha,H}\), the definition of \(r_{n,\alpha,H}\) gives
\[
\|\Gamma_e-\Sigma_e\|_{\mathrm{op}}\leq r
\qquad(e\in S).
\tag{9}
\]
Using the candidate equation, (2), (9), submultiplicativity, and \(\|X\|_{\mathrm F}\leq\sqrt p\|X\|_{\mathrm{op}}\), one obtains for each \(e\in S\)
\[
\begin{aligned}
\left\|A\Sigma_eA^{\mathsf T}
-\Psi-\operatorname{diag}(t_e)\right\|_{\mathrm F}
&=\left\|A(\Sigma_e-\Gamma_e)A^{\mathsf T}\right\|_{\mathrm F}\\
&\leq\sqrt p\,L_0^2r
\leq\mu_0/2.
\end{aligned}
\tag{10}
\]
If this tuple were in \(\mathfrak T_{\mathrm{far}}\), then (5) would force the maximum of the left-hand side in (10) to be at least \(\mu_0\) when the far family is nonempty; membership is impossible by definition when it is empty. Thus, with
\[
Q=AD^{-1}=I_p+R,
\qquad u=\|R\|_{\mathrm F},
\]
one has
\[
u<\rho_0,
\qquad K_0\overline M u<\frac12.
\tag{11}
\]

\textbf{Local affine-minor inversion.}
Fix \(e_0\in S\) and write \(x_e=s_e-s_{e_0}\). For each \(i<j\), affine separation supplies \(a,b,d\in S\) such that the determinant of the three points \((s_{ri},s_{rj})\), \(r\in\{a,b,d\}\), has absolute value at least \(\gamma_0\). Translation by \(-(s_{e_0i},s_{e_0j})\) preserves this determinant. Moreover,
\[
\det(\widetilde x_a-\widetilde x_b,
\widetilde x_d-\widetilde x_b)
=\det(\widetilde x_a,\widetilde x_d)
+\det(\widetilde x_d,\widetilde x_b)
+\det(\widetilde x_b,\widetilde x_a),
\]
where \(\widetilde x_e=(x_{ei},x_{ej})\). The triangle inequality therefore implies that two labels \(e,f\in S\) obey
\[
|x_{ei}x_{fj}-x_{ej}x_{fi}|\geq\gamma_0/3.
\tag{12}
\]
Let \(E_e=\Gamma_e-\Sigma_e\). The true difference equation is
\[
D(\Sigma_e-\Sigma_{e_0})D^{\mathsf T}=\operatorname{diag}(x_e),
\]
and the candidate difference equation is diagonal after adding \(A(E_e-E_{e_0})A^{\mathsf T}\). Since \(A=QD\), the \((i,j)\) off-diagonal entry, for \(i\neq j\), gives exactly
\[
R_{ij}x_{ej}+R_{ji}x_{ei}
=-\sum_{k=1}^pR_{ik}R_{jk}x_{ek}
-[A(E_e-E_{e_0})A^{\mathsf T}]_{ij}.
\tag{13}
\]
The true scale bound and nonnegativity imply \(0\leq s_{ek}\leq\overline M\), hence \(|x_{ek}|\leq\overline M\). Equation (9) gives \(\|E_e-E_{e_0}\|_{\mathrm{op}}\leq2r\), and (2) gives \(\|A\|_{\mathrm{op}}\leq L_0\). Cauchy--Schwarz also gives \(\sum_k|R_{ik}R_{jk}|\leq u^2\). Therefore (13) implies
\[
|R_{ij}x_{ej}+R_{ji}x_{ei}|
\leq\overline M u^2+2rL_0^2.
\tag{14}
\]
Apply (14) at the two labels from (12). The coefficient determinant of this two-by-two system has absolute value at least \(\gamma_0/3\), while every coefficient has absolute value at most \(\overline M\). Cramer's rule yields, for every ordered pair \(i\neq j\),
\[
|R_{ij}|
\leq\frac{6\overline M}{\gamma_0}
\left(\overline M u^2+2rL_0^2\right).
\tag{15}
\]
The unit diagonals of \(A\) and \(D\) imply
\[
0=(A-D)_{ii}=(RD)_{ii}
=R_{ii}+\sum_{k\neq i}R_{ik}D_{ki}.
\tag{16}
\]
For each \(i\), the squared norm of the \(i\)-th column of \(D\) is at most \(L_0^2\). Square (15), sum over the \(p(p-1)\) ordered off-diagonal entries, apply Cauchy--Schwarz to (16), and take square roots. By the definition of \(K_0\), this gives
\[
u\leq K_0\left(\overline M u^2+2rL_0^2\right).
\tag{17}
\]
By (11), \(K_0\overline M u^2<u/2\). Moving this term to the left in (17) gives
\[
u\leq4K_0L_0^2r.
\tag{18}
\]
It follows that
\[
\|A-D\|_{\mathrm{op}}
=\|RD\|_{\mathrm{op}}
\leq uL_0
\leq4K_0L_0^3r
\leq C_0r.
\tag{19}
\]
Taking the supremum over \(A\in\mathcal U^{(n)}_\alpha\) proves the outer-radius bound. On \(G_{n,\alpha,H}\), \(\text{thm:confidence-union-coverage}\) gives \(D\in\mathcal U^{(n)}_\alpha\), so the union is nonempty. For any \(A_1,A_2\in\mathcal U^{(n)}_\alpha\), the triangle inequality and (19) yield
\[
\|A_1-A_2\|_{\mathrm{op}}
\leq\|A_1-D\|_{\mathrm{op}}+\|A_2-D\|_{\mathrm{op}}
\leq2C_0r.
\]
This proves the diameter bound. The same \(r_0\) and \(C_0\) apply throughout the stated fixed-margin class, so \(r_{n,\alpha,H}\to0\) along covered outcomes implies convergence of both set radii to zero. Every step used only positive semidefiniteness, bounded scale, affine separation, cycle slack, and structural conditioning; no covariance inverse or lower covariance eigenvalue entered the argument.